% ============================================================
% 第 8 章 交流绕组（教材第三篇，印刷页 78--89）
% 节名：第一节 同步发电机和感应电动机工作原理概述；第二节 三相单层绕组；第三节 三相双层绕组
% ============================================================
\chapter{交流绕组}

工业用交流电机主要分为两大类：\textbf{同步电机}与\textbf{感应电机}。两者工作原理和运行特点不同，但有一批共性问题：定子绕组都是三相交流绕组，绕组中感应电动势与磁动势的分析计算方法相同，发热、冷却与温升限值的规律也相同。本篇（第 8—11 章）解决这些共性问题，第 12 章起再分别研究两类电机的运行性能。

\section{同步发电机和感应电动机工作原理概述}

\subsection{同步发电机的基本工作原理}
三相同步发电机由\textbf{定子}、\textbf{转子}和定转子之间的工作气隙组成。定子是一个圆筒形铁芯，靠近铁芯内表面均匀开槽，槽里嵌放导体；把这些导体按一定规律连接起来构成\textbf{定子绕组}（也称电枢绕组）。最简单的三相定子绕组每相仅有一个线圈，首末端分别用 $A$—$X$、$B$—$Y$、$C$—$Z$ 表示，它们在空间位置上彼此相差 $120^\circ$。

圆筒形铁芯的中间是可旋转的转子，转子上装有\textbf{主磁极}，磁极铁芯上套有励磁绕组。给励磁绕组通入直流电流，磁极就表现出极性（$N$、$S$）。在 $N$ 极处磁力线由转子磁极表面流出进入定子铁芯，在 $S$ 极下磁力线由定子铁芯流出进入转子磁极，形成一对极的磁场。

当转子由原动机拖动以恒定转速 $n_1$（\textbf{同步转速}）旋转时，定子绕组切割转子磁场感应交流电动势，感应电动势方向由右手定则确定。转子每旋转一周，感应电动势交变一次；转速单位为 r/min，故
\begin{equation}
f=\frac{pn_1}{60}\ \text{(Hz)} .
\end{equation}
又因三相绕组在空间上互差 $120^\circ$，转子磁场被 $A$ 相绕组切割后再转过 $120^\circ$ 被 $B$ 相切割，故 $B$ 相电动势滞后 $A$ 相 $120^\circ$，$C$ 相又滞后 $B$ 相 $120^\circ$，得到三相对称感应电动势。

同步发电机三相绕组通常接成星形，加上三相对称负载后，在三相对称电动势作用下向负载输出电能；与此同时，三相绕组中的电流与磁场相互作用，产生与原动机拖动转矩相反的\textbf{电磁转矩}，克服原动机的机械能，把机械能转换成电能，形成机电能量转换。

\subsection{感应电动机的基本工作原理}
感应电动机的定子结构与同步发电机相同，由铁芯和交流绕组构成。转子也由铁芯和绕组构成，但结构与同步发电机不同：以常用的三相笼型感应电动机为例，靠近转子铁芯外表面均匀开槽，槽里嵌放\textbf{导条}（铜条或铸铝），导条两端用\textbf{端环}连接，形成笼型绕组。

\subsubsection{三相对称电流产生旋转磁场}
约定：当 $A$ 相电流为正值时，电流从 $A$ 导体流入纸面、从 $X$ 导体流出纸面；$B$、$C$ 相类似。

某时刻 $A$ 相电流为最大值 $I_m$、$B$、$C$ 相电流均为 $-I_m/2$。此时 $A$ 相电流从 $A$ 端流入、$X$ 端流出；$B$、$C$ 相电流分别从 $Y$、$Z$ 端流入、从 $B$、$C$ 端流出。由右手定则可判断，三相绕组电流产生的合成磁动势处于\textbf{$A$ 相绕组的轴线}上。此时 $A$ 相电流最大，对合成磁动势贡献最大；$B$、$C$ 相瞬时电流只有 $I_m/2$，且其作用方向与合成磁动势方向成 $60^\circ$ 夹角，故两相对合成磁动势的贡献分别为 $\tfrac12F_{\phi1}\cos60^\circ$。于是合成磁动势
\begin{equation}
F=\frac32F_{\phi1},
\end{equation}
即合成磁动势为单相绕组产生磁动势最大值的 $3/2$ 倍。

再过 $30^\circ$（时间上），合成磁动势位置与前一时刻相比\textbf{旋转了 $30^\circ$}；此时合成磁动势由 $A$ 相与 $B$（$C$）相电流共同产生，两者作用方向与合成磁动势的夹角都是 $30^\circ$，合成磁动势仍为 $\tfrac32F_{\phi1}$。既然 $A$ 相电流最大时磁动势最大值在 $A$ 相轴线上，那么 $B$ 相电流达最大时合成磁动势必在 $B$ 相轴线上，$C$ 相电流最大时在 $C$ 相轴线上。以此类推：\textbf{电流变化一个周期（$2\pi$），磁动势在空间旋转一周，且幅值始终保持为 $\tfrac32F_{\phi1}$}。该磁动势产生两极磁场，其转速
\begin{equation}
n_1=\frac{60f}{p}\ \text{(r/min)} .
\end{equation}
这就是三相对称交流绕组通以三相对称交流电流产生旋转磁场的基本概念。

\subsubsection{转子转动的物理过程}
把定子旋转磁场用旋转磁极表示，设其逆时针同步旋转、转速 $n_1$。当转子静止（$n=0$）或转速 $n<n_1$ 时，转子导条与旋转磁场有相对运动，导条中感应电动势。因转子速度低于 $n_1$，导条切割磁力线的运动方向为顺时针，由右手定则可判断感应电动势方向；导条被端环短接，若不计导条电抗影响，导条中电流与感应电动势同相位。载流导体在磁场中受电磁力，由左手定则可判断方向，电磁力与转子半径之积即\textbf{电磁转矩}。转子上的电磁转矩为驱动力矩，若能克服负载和摩擦阻力矩，电机就以转速 $n$ 旋转。

因为转子电路中没有外接电源，完全依靠转子与旋转磁场之间的相对运动感应电动势、电流并产生电磁转矩，故称\textbf{感应电动机}；又因为带动机械负载时转速总要低于同步转速 $n_1$，故又称\textbf{异步电动机}。

反之，定子绕组 $A$—$X$、$B$—$Y$、$C$—$Z$ 三相对称（即三相绕组导体沿定子内表面按 $AZBXCY$ 规律排列），当转子以同步速旋转时，定子绕组中将产生三相对称感应电动势。

\subsection{多极电机}
实际交流电机制成多极。多极电机可视作若干个两极电机拼合而成：四极电机可看作两个两极电机拼合，各占半个圆周（机械角 $180^\circ$）。转子转过一对极（四极电机的半周），定子绕组感应电动势交变一次；转子转过一周，电动势交变两次。一般地，对 $p$ 对极电机，转子转过一周感应电动势交变 $p$ 次，故
\begin{equation}
f=\frac{pn}{60}.
\end{equation}

感应电动机中，两极电机电流交变一次、磁场旋转一周（转过一对极）；$p$ 对极电机可视作 $p$ 个两极电机拼合，电流交变一次磁场仍转一对极（$1/p$ 圆周），故 $p$ 对极电机旋转磁场转速为
\begin{equation}
n_1=\frac{60f}{p}.
\end{equation}

\keyformula{n_1=\frac{60f}{p}\ \text{(r/min)},\qquad f=\frac{pn}{60}}

两式实际上是同一公式的两种写法：对同步发电机，转子旋转一周定子电动势变化 $p$ 次；对感应电动机，定子电流交变一次磁场旋转 $1/p$ 周。即电气量随时间变化 $360^\circ$，则磁场或转子在空间旋转 $360^\circ/p$。若用电气量变化的角度（\textbf{电角度}）表示空间机械角度，则
\begin{equation}
\text{电角度}=p\times\text{空间机械角度} .
\end{equation}

引入空间电角度的好处：无论一对极还是 $p$ 对极电机，用电角度表示空间机械角度后，该角度与对应的感应电动势（或电流）随时间变化的角度一致。另外，相邻两相绕组空间位置上的机械角差为 $120^\circ/p$，空间电角差为 $120^\circ$，恰好等于相邻两绕组中感应电动势的时间相位差——\textbf{时间相位差与空间电角度差相一致}，这是交流绕组分析的总纲。

\corepoint{三句话背下来，第 8 章就通了一半：① $n_1=60f/p$ 描述"旋转磁场多快"；② $f=pn/60$ 描述"感应电动势多快"；③ 电角度 $=p\times$ 机械角度描述"角度怎么换算"。凡是遇到"绕组空间位置差多少度、电动势相位差多少度"，一律先换到电角度再算，千万不要把机械角度当电角度用。}

\section{三相单层绕组}

\subsection{从集中绕组到分布绕组}
图 8-1、8-5 中的两极、四极三相定子绕组都属于\textbf{集中绕组}，实际电机中的绕组比这复杂。集中绕组有很多缺点：电气性能差、散热不好，而且浪费铜线、铁芯等有效材料。实际电机大多采用\textbf{分布绕组}，它由集中绕组演变而来：把原来集中在一个大槽中的导体分散放置在 $q$ 个小槽中，$q$ 个小槽对应的圆周角度称为\textbf{相带}。所有大槽都这样处理，就得到 $Q=2pmq$ 个小槽构成的分布绕组；每小槽的导体都属于同一线圈，故称\textbf{单层绕组}。属于同一相的 $q$ 个槽分布在 $60^\circ$ 空间电角度范围内，故称 \textbf{$60^\circ$ 相带绕组}。

分布绕组各相带内导体的连接应与集中时相同，即每相带 $q$ 个导体依次串联，串联的先后次序可以改变。为节省铜线、制造方便和端部布置合理，按线圈形状和端部连接方法，单层绕组可分为\textbf{同心式、链式和交叉式}。

\subsection{几个基本量}
\begin{enumerate}
  \item \textbf{极距} $\tau$：对应于一个极（即 $180^\circ$ 空间电角度）的定子内表面圆弧长度，可用槽数或圆弧长度度量：
  \begin{equation}
  \tau=\frac{Q}{2p}\ \text{(槽)}\qquad\text{或}\qquad \tau=\frac{\pi D}{2p},
  \end{equation}
  其中 $Q$ 为定子槽数，$D$ 为定子铁芯内径。
  \item \textbf{节距} $y$：一个线圈的两个线圈边之间的圆弧长度，可用槽数或圆弧长度度量。$y=\tau$ 为整距，$y<\tau$ 为短距。
  \item \textbf{每极每相槽数} $q$：每个极下每相所占的槽数，即由集中绕组过渡到分布绕组时的小槽数：
  \begin{equation}
  q=\frac{Q}{2pm},
  \end{equation}
  其中 $m$ 为相数，三相电机 $m=3$。
  \item \textbf{槽距角} $\alpha$：相邻两槽之间的电角度，
  \begin{equation}
  \alpha=\frac{p\times360^\circ}{Q}.
  \end{equation}
\end{enumerate}

\subsection{同心式绕组}
同心式绕组由不等节距的同心线圈构成。例如三相二极 24 槽的定子，每极每相槽数
\begin{equation}
q=\frac{Q}{2pm}=\frac{24}{2\times1\times3}=4 .
\end{equation}
按 $60^\circ$ 相带原则划分为 6 个相带，各相带内的槽号按顺序排列（每相带 4 槽：$23,24,1,2$ / $3,4,5,6$ / …）。

由于相带 $A$ 与相带 $X$ 处于相异磁极下，其电流方向相反：设导体 $23,24,1,2$ 中电流流入纸面，则导体 $11,12,13,14$ 中电流流出纸面。把 $1$ 和 $12$ 相连构成一个大线圈，$2$ 和 $11$ 相连构成一个小线圈，这一大一小构成一个\textbf{同心式线圈组}；再把 $13$ 与 $24$ 相连、$14$ 与 $23$ 相连，构成另一个同心式线圈组；然后顺着电流流向把两个线圈组连接起来（尾端与尾端相接、首首引出），作为 $A$、$X$ 端，即构成 $A$ 相绕组。其余两相同理。

\subsection{链式与交叉式绕组}
\begin{itemize}
  \item \textbf{链式绕组}：线圈节距相同、端部如链环相扣，各线圈依次串联，端部较短、省铜线。
  \item \textbf{交叉式绕组}：由大小两种节距的线圈交叉组合，端部排列更紧凑，多用于 $q$ 为奇数的场合。
\end{itemize}

单层绕组的主要优点是：\textbf{槽内无须层间绝缘、槽利用率高、嵌线方便}，多用于散下线的 10\,kW 以下小功率交流电机；主要缺点是\textbf{电气性能不如双层绕组}——不能通过选择节距来削弱谐波电动势与磁动势。

\section{三相双层绕组}

双层绕组的每个槽分上、下两层，每槽内放置两个线圈边（一个上层边、一个下层边）；线圈数等于槽数，一个线圈的一个边放在某槽上层，另一个边放在相隔节距 $y$ 的另一槽下层。

\subsection{结构特点}
\begin{enumerate}
  \item \textbf{可选择节距}：双层绕组可以做成短距（通常 $y_1\approx\tfrac56\tau$ 或 $\tfrac{5}{6}\tau$ 附近），从而\textbf{同时削弱 5 次和 7 次谐波}电动势与磁动势，改善波形（第 9 章第五节将定量说明）。
  \item \textbf{端部排列整齐}：所有线圈形状相同，端部易于布置、便于制造。
  \item \textbf{分叠绕组与波绕组}：叠绕组是相邻线圈依次串联、后一线圈紧叠在前一线圈上；波绕组是线圈沿圆周"波浪式"前进，多用于极数较多、支路数要求特殊的场合。
\end{enumerate}

在绕组展开图中，上层边用实线表示、下层边用虚线表示，每一个线圈都是一半实线、一半虚线。按 $60^\circ$ 相带划分后，各相带槽号的分配与单层绕组相同，只是每个槽号对应上、下两个线圈边。

\subsection{两种绕组对比}
\begin{center}
\begin{tabular}{@{}p{2.2cm}p{3.6cm}p{3.9cm}p{3.0cm}@{}}
\toprule
项目 & 单层绕组 & 双层绕组 & 说明 \\
\midrule
每槽线圈边 & 1 个 & 2 个（上、下层） & 双层槽利用率略低 \\
线圈数 & $Q/2$ & $Q$ & 双层线圈数等于槽数 \\
节距 & 一般为整距 & 可选短距（$\approx\tfrac56\tau$） & 短距可削弱 5、7 次谐波 \\
端部 & 同心式/链式/交叉式 & 叠绕组/波绕组 & 双层端部整齐 \\
电气性能 & 较差 & 好 & 大、中型电机多用双层 \\
典型用途 & 10\,kW 以下小功率电机 & 大、中型交流电机 & — \\
\bottomrule
\end{tabular}
\end{center}

\corepoint{绕组这一章的所有计算都归结为四个量：$p$（极对数）、$Q$（槽数）、$m$（相数）、$y$（节距）。由它们派生出 $\tau=Q/(2p)$、$q=Q/(2pm)$、$\alpha=360^\circ p/Q$。做题第一步永远是先算这四个派生量并写清楚，再画槽号分配表；槽号分配错了，后面电动势、磁动势全错。另外牢记：单层绕组\textbf{不能}用短距削弱谐波，双层才可以——这是选择绕组型式的第一条理由。}
